% !TeX root = ../../Book3.tex
% 纲要标题：### A.4 理想类群与理想范数的 Euler 乘积

\section{理想类群与理想范数的 Euler 乘积}
\label{b3:sec:A04}

% 纲要细目（与计划纲要同步）：
% * 类域论与全局解析工具不纳入本卷核心证明链
%
% <!-- 短节：不设小节 -->


理想分解修复了元素分解失效的问题，但仍留下一个具体问题：一个理想何时可以由单个元素生成？把主理想与一般理想的差别组织成群，便得到继续学习代数数论所需的一个基本对象。

\ProseParagraphBreak
沿用\cref{b3:def:fractional-invertible-ideal,b3:def:ideal-class-picard}。在数域 \(K\) 中，令 \(R=\mathcal O_K\)；非零分式理想可以等价地看作 \(K\) 的非零有限生成 \(R\)-子模，或可用一个非零 \(d\in R\) 清除分母的非零子模。乘法由有限和 \(\sum a_ib_i\) 定义；逆元为
\[
I^{-1}=\{x\in K:xI\subseteq R\}.
\]
\cref{b3:cor:dedekind-fractional-group}已经证明它们构成交换群，且每个分式理想唯一分解成素理想的整数次幂。理想类群是商群
\[
\operatorname{Cl}(R)=
\{\text{非零分式理想}\}/\{aR:a\in K^\times\}.
\]
同一理想类中的理想相差一个非零标量。群平凡恰好表示每个非零理想都主生成，见\cref{b3:prop:dedekind-picard-class}。

\begin{example}\label{b3:exm:nonprincipal-ideal-minus-six}
在 \(R=\mathbb Z[\sqrt{-6}]\) 中，令
\(\mathfrak p=(2,\sqrt{-6})\)。二次整基公式说明 \(R\) 确为整数环，商环为 \(\mathbb F_2\)，故 \(N(\mathfrak p)=2\)。若 \(\mathfrak p=(a+b\sqrt{-6})\)，主理想的指数等于乘法矩阵行列式的绝对值，因此
\[
2=|N_{K/\mathbb Q}(a+b\sqrt{-6})|=a^2+6b^2,
\]
这没有整数解。所以 \(\mathfrak p\) 非主。

\ProseParagraphBreak
另一方面，\(\mathfrak p^2\) 由
\(4,\,2\sqrt{-6},\,-6\) 生成，均属于 \(2R\)，而 \(-6+2\cdot4=2\)。故 \(\mathfrak p^2=2R\)。类群中 \([\mathfrak p]\) 的阶恰为 \(2\)。这个计算与第10章的 \(\mathbb Q(\sqrt{-5})\) 例子使用同一方法：先以剩余域确定理想范数，再用一个整数二次型排除主生成。它只证明存在这样的非平凡理想类，并没有枚举整个类群。
\end{example}

理想类群把“差多少才成为主理想”变成一个群问题。进一步的类域论把理想类及带同余条件的类似对象，与数域的交换 Galois 扩张联系起来。完整的陈述需要定义射线类群、局部与全局互反映射，并逐项说明无穷素点及分歧的处理；本附录不把这些对象压缩成一个没有准备的对应定理。当前已经建立的理想分解、赋值和完备化，正是理解那些定义所需的算术基础。

\ProseParagraphBreak
另一种研究方法把全部理想范数放进一个级数。它从唯一分解出发，与素数分布和理想类的计数问题发生联系。下面只证明绝对收敛区域中的乘积公式，不使用解析延拓或类数公式。

\begin{proposition}\label{b3:prop:dedekind-zeta-euler}
设 \(K/\mathbb Q\) 的次数为 \(n\)，\(\mathcal O_K\) 为其代数整数环，并记非零理想 \(I\subseteq\mathcal O_K\) 的范数为 \(N(I)=|\mathcal O_K/I|\)。对实数 \(s>1\)，级数
\[
\zeta_K(s)=\sum_{0\ne I\subseteq\mathcal O_K}N(I)^{-s}
\]
收敛，并满足
\[
\zeta_K(s)=\prod_{\mathfrak p}
\bigl(1-N(\mathfrak p)^{-s}\bigr)^{-1},
\]
其中乘积遍历非零素理想。
\end{proposition}
\begin{proof}
理想分解、中国剩余定理及素理想幂的范数公式给出范数的乘法性。对于任意有限组素理想，展开有限个几何级数，所得和恰好遍历只由这些素理想组成的理想，且所有项非负。

一个有理素数 \(p\) 上方至多有 \(n\) 个素理想，各范数至少为 \(p\)。因此任意有限子乘积都不超过
\[
\prod_p(1-p^{-s})^{-n}=\zeta(s)^n<\infty.
\]
这里最后一个等式由整数唯一分解同样得到，普通级数
\(\sum_{m\ge1}m^{-s}\) 在 \(s>1\) 时由积分比较收敛。有限子乘积单调有界；每个非零理想只含有限多个素因子，故取极限恰得到所列级数与乘积。
\end{proof}

这个公式的代数内容已经全部体现在理想唯一分解中。若要研究 \(s=1\) 附近的行为，则需要额外的分析工具；若要由它计算类数，还需要单位群及数域的实、复嵌入。后续学习可以分别从理想类、局部域或这个级数出发，但这三方面的深入定理均不作为本卷正文的证明依据。
